\documentclass{scrartcl}
\usepackage{graphicx}  % required for inserting images

% for cyrillic and Bulgarian language
\usepackage[utf8]{inputenc}
\usepackage[T2A]{fontenc}
\usepackage[bulgarian]{babel}

% math packages
\usepackage{amsmath}
\usepackage{float}
\usepackage{amsfonts}
\usepackage{amssymb}
\usepackage{amsthm}
\usepackage{cancel}

\usepackage{ragged2e}  % adds FlushLeft

\usepackage{listings}
\usepackage{xcolor}

\definecolor{codegreen}{rgb}{0,0.6,0}
\definecolor{codegray}{rgb}{0.5,0.5,0.5}
\definecolor{codepurple}{rgb}{0.58,0,0.82}
\definecolor{backcolour}{rgb}{0.95,0.95,0.92}

\lstdefinestyle{mystyle}{
    language=Octave,
    backgroundcolor=\color{backcolour},   
    commentstyle=\color{codegreen},
    keywordstyle=\color{magenta},
    numberstyle=\tiny\color{codegray},
    stringstyle=\color{codepurple},
    basicstyle=\ttfamily\footnotesize,
    breakatwhitespace=false,         
    breaklines=true,
    captionpos=b,
    keepspaces=false,                 
    numbersep=5pt,
    showspaces=false,
    showstringspaces=false,
    showtabs=false,
    tabsize=2,
    columns=fullflexible,
}

\lstset{style=mystyle}

\usepackage[colorlinks=true, linkcolor=blue, urlcolor=cyan]{hyperref}

\newtheorem{theorem}{Теорема}
\newtheorem{definition}{Дефиниция}

\newtheorem{lemma}{Лема}
\newtheorem{example}{Пример}
\newtheorem{remark}{Забележка}

\title{ДИС - Теорема на Ферма. Теорема за средни стойности. Формула на Тейлър}
\author{Ивайло Андреев + gemini 3 Flash}
\date{16 май 2026}

\begin{document}
\maketitle

\tableofcontents

\newpage

\begin{FlushLeft}

\section{Локални екстремуми и Теорема на Ферма}

Математическият анализ на функциите често изисква изследване на техните най-големи и най-малки стойности в локален смисъл. Първата стъпка към това е дефинирането на локален екстремум.

\begin{definition}[Локален екстремум]
Нека $f: D \to \mathbb{R}$ е функция, дефинирана върху множеството $D \subset \mathbb{R}$, и $x_0 \in D$ е вътрешна точка за $D$.
\begin{enumerate}
    \item Казваме, че $f$ има \textbf{локален минимум} в точката $x_0$, ако съществува $\delta > 0$ такова, че околността $(x_0 - \delta, x_0 + \delta) \subset D$ и за всяко $x \in (x_0 - \delta, x_0 + \delta)$ е изпълнено:
    \[ f(x_0) \le f(x) \]
    \item Казваме, че $f$ има \textbf{локален максимум} в точката $x_0$, ако съществува $\delta > 0$ такова, че околността $(x_0 - \delta, x_0 + \delta) \subset D$ и за всяко $x \in (x_0 - \delta, x_0 + \delta)$ е изпълнено:
    \[ f(x_0) \ge f(x) \]
\end{enumerate}
Локалните минимуми и максимуми се обединяват под общото наименование \textbf{локални екстремуми}.
\end{definition}

\begin{theorem}[Теорема на Ферма - Необходимо условие за локален екстремум]
Нека функцията $f$ е дефинирана в околност на точката $x_0$ и има локален екстремум в $x_0$. Ако $f$ е диференцируема в точката $x_0$, то:
\[ f'(x_0) = 0 \]
\end{theorem}

\begin{proof}
Ще проведем доказателството за случая, когато $f$ има \textbf{локален максимум} в точката $x_0$ (случаят за локален минимум се доказва напълно аналогично).

По дефиниция, щом $f$ има локален максимум в $x_0$, съществува околност $(x_0 - \delta, x_0 + \delta)$, в която:
\[ f(x) \le f(x_0) \quad \Rightarrow \quad f(x) - f(x_0) \le 0 \]
за всяко $x$ от тази околност.

По условие функцията е диференцируема в $x_0$, което означава, че съществува границата на диференчното частно и тя е равна на производната $f'(x_0)$:
\[ f'(x_0) = \lim_{h \to 0} \frac{f(x_0 + h) - f(x_0)}{h} \]
Тъй като тази граница съществува, левите и десните едностранни граници също съществуват и са равни на $f'(x_0)$. Нека изследваме знака на диференчното частно за $h \in (-\delta, \delta) \setminus \{0\}$:

\begin{enumerate}
    \item \textbf{Случай 1 (Дясна граница):} Нека $h > 0$. Тогава:
    \[ \frac{f(x_0 + h) - f(x_0)}{h} = \frac{\le 0}{> 0} \le 0 \]
    Преминавайки към граница при $h \to +0$, от теоремата за запазване на знака при граничен преход получаваме:
    \[ f'(x_0) = \lim_{h \to +0} \frac{f(x_0 + h) - f(x_0)}{h} \le 0 \]
    
    \item \textbf{Случай 2 (Лява граница):} Нека $h < 0$. Тогава:
    \[ \frac{f(x_0 + h) - f(x_0)}{h} = \frac{\le 0}{< 0} \ge 0 \]
    Преминавайки към граница при $h \to -0$, получаваме:
    \[ f'(x_0) = \lim_{h \to -0} \frac{f(x_0 + h) - f(x_0)}{h} \ge 0 \]
\end{enumerate}

Тъй като $f'(x_0)$ е едно и също число (поради диференцируемостта), от двете получени неравенства $f'(x_0) \le 0$ и $f'(x_0) \ge 0$ по правилата за сравняване на реални числа следва единствената възможност:
\[ f'(x_0) = 0 \]
С това теоремата е доказана.
\end{proof}

\begin{remark}[Геометричен смисъл]
Геометрично теоремата на Ферма означава, че в точката на локален екстремум допирателната към графиката на функцията $y = f(x)$ е хоризонтална линия, т.е. успоредна на абсцисната ос $Ox$.
\end{remark}

\newpage

\section{Теореми за средните стойности}

Теоремите на Рол, Лагранж и Коши съставляват ядрото на диференциалното смятане. Те установяват връзка между стойностите на функцията в краищата на даден затворен интервал и стойностите на нейната производна във вътрешни точки на интервала.

За трите теореми се изисква една фундаментална помощна теорема, която приемаме без доказателство съгласно изискванията на конспекта:

\begin{theorem}[Втора теорема на Вайерщрас]
Всяка функция, която е непрекъсната в краен и затворен интервал $[a, b]$, достига в него своята точна горна граница (максимум) $M$ и своята точна долна граница (минимум) $m$. Тоест, съществуват точки $x_1, x_2 \in [a, b]$, такива че $f(x_1) = m$ и $f(x_2) = M$.
\end{theorem}

\subsection{Теорема на Рол}

\begin{theorem}[Теорема на Рол]
Нека функцията $f(x)$ отговаря на следните три условия:
\begin{enumerate}
    \item $f(x)$ е непрекъсната в затворения интервал $[a, b]$;
    \item $f(x)$ притежава производна (диференцируема е) поне в отворения интервал $(a, b)$;
    \item Стойностите в краищата на интервала са равни: $f(a) = f(b)$.
\end{enumerate}
Тогава съществува поне една точка $c \in (a, b)$, такава че:
\[ f'(c) = 0 \]
\end{theorem}

\begin{proof}
Тъй като $f(x)$ е непрекъсната в затворения интервал $[a, b]$, съгласно Теоремата на Вайерщрас тя достига своя минимум $m$ и своя максимум $M$ в този интервал. Нека:
\[ m = \min_{x \in [a, b]} f(x) = f(x_1), \quad M = \max_{x \in [a, b]} f(x) = f(x_2) \quad (x_1, x_2 \in [a, b]) \]

Разглеждаме два основни случая за стойностите на $m$ и $M$:

\begin{enumerate}
    \item \textbf{Случай 1:} $m = M$.
    Ако минимумът е равен на максимума, то функцията е константа в целия интервал: $f(x) = c = \text{const}$ за всяко $x \in [a, b]$. Тогава нейната производна е нула във всяка точка от интервала: $f'(x) = 0$. В този случай всяка точка от отворения интервал $(a, b)$ може да бъде избрана за точка $c$.
    
    \item \textbf{Случай 2:} $m < M$.
    Тъй като по условие $f(a) = f(b)$, не е възможно и двете стойности $m$ и $M$ да се достигат едновременно в краищата на интервала. Следователно, поне една от двете стойности (или минимумът $m$, или максимумът $M$) се достига във вътрешна точка $c \in (a, b)$.
    
    Тъй като $c$ е вътрешна точка на интервала, тя се явява точка на локален екстремум (локален минимум или локален максимум) за функцията $f$. Освен това, по условие функцията е диференцируема в отворения интервал $(a, b)$, следователно тя е диференцируема и в точката $c$. 
    
    Прилагаме директно \textbf{Теоремата на Ферма} за точката $c$, откъдето следва:
    \[ f'(c) = 0 \]
\end{enumerate}
С това теоремата на Рол е доказана.
\end{proof}

\subsection{Теорема на Лагранж}

Теоремата на Лагранж (често наричана "Формула за крайните нараствания") премахва ограничението $f(a) = f(b)$ от теоремата на Рол.

\begin{theorem}[Теорема на Лагранж]
Нека функцията $f(x)$ е непрекъсната в затворения интервал $[a, b]$ и диференцируема в отворения интервал $(a, b)$. Тогава съществува поне една точка $c \in (a, b)$, такава че:
\[ f(b) - f(a) = f'(c)(b - a) \]
\end{theorem}

\begin{proof}
За доказателството ще използваме метода на помощната функция, като я конструираме така, че за нея да са изпълнени условията на теоремата на Рол.

Нека дефинираме числото $k$, представляващо наклона на секущата през краищата на графиката:
\[ k = \frac{f(b) - f(a)}{b - a} \]
Конструираме помощната функция $F(x)$ по следния начин (както е в официалните лекции):
\[ F(x) = f(x) - k \cdot x = f(x) - \frac{f(b) - f(a)}{b - a} \cdot x \]

Проверяваме условията на теоремата на Рол за функцията $F(x)$ в интервала $[a, b]$:
\begin{enumerate}
    \item $F(x)$ е непрекъсната в $[a, b]$, защото е разлика на непрекъснатата функция $f(x)$ и линейната функция $k \cdot x$.
    \item $F(x)$ е диференцируема в $(a, b)$, тъй като е разлика на диференцируеми функции. Нейната производна е:
    \[ F'(x) = f'(x) - k = f'(x) - \frac{f(b) - f(a)}{b - a} \]
    \item Проверяваме стойностите в краищата на интервала:
    \[ F(a) = f(a) - \frac{f(b) - f(a)}{b - a} \cdot a = \frac{f(a)b - f(a)a - f(b)a + f(a)a}{b - a} = \frac{f(a)b - f(b)a}{b - a} \]
    \[ F(b) = f(b) - \frac{f(b) - f(a)}{b - a} \cdot b = \frac{f(b)b - f(a)b - f(b)b + f(a)b}{b - a} = \frac{f(a)b - f(b)a}{b - a} \]
    Виждаме, че $F(a) = F(b)$.
\end{enumerate}

Тъй като всички условия на теоремата на Рол са изпълнени за $F(x)$, то съществува точка $c \in (a, b)$, такава че $F'(c) = 0$. Заместваме $c$ в израза за производната:
\[ F'(c) = f'(c) - \frac{f(b) - f(a)}{b - a} = 0 \quad \Rightarrow \quad f'(c) = \frac{f(b) - f(a)}{b - a} \]
Умножавайки двете страни на равенството по $(b - a)$, получаваме търсената формула:
\[ f(b) - f(a) = f'(c)(b - a) \]
Теоремата е доказана.
\end{proof}

\subsection{Теорема на Коши}

Теорема на Коши е най-общата от трите теореми за средните стойности и разглежда две функции едновременно.

\begin{theorem}[Теорема на Коши]
Нека функциите $f(x)$ и $g(x)$ са непрекъснати в затворения интервал $[a, b]$ и диференцируеми в отворения интервал $(a, b)$, като $g'(x) \neq 0$ за всяко $x \in (a, b)$. Тогава съществува точка $c \in (a, b)$, такава че:
\[ \frac{f(b) - f(a)}{g(b) - g(a)} = \frac{f'(c)}{g'(c)} \]
\end{theorem}

\begin{proof}
Доказателството се състои от две задължителни логически стъпки съгласно изискванията на конспекта.

\textbf{Стъпка 1: Доказателство, че знаменателят $g(b) - g(a) \neq 0$.}
Ще докажем това твърдение с допускане на противното. Да допуснем, че $g(b) - g(a) = 0$, което означава $g(a) = g(b)$.
За функцията $g(x)$ са изпълнени всички условия на теоремата на Рол в интервала $[a, b]$ (тя е непрекъсната в $[a, b]$, диференцируема в $(a, b)$ и $g(a) = g(b)$). Тогава, по теоремата на Рол, трябва да съществува точка $\xi \in (a, b)$, за която $g'(\xi) = 0$.
Това обаче директно противоречи на условието на теоремата, че $g'(x) \neq 0$ за всяко $x \in (a, b)$. Следователно допускането е невярно и:
\[ g(b) - g(a) \neq 0 \]

\textbf{Стъпка 2: Използване на помощна функция.}
Дефинираме помощната функция $H(x)$ (съобразно структурата от лекциите и разработките):
\[ H(x) = f(x) - \lambda \cdot g(x), \quad \text{където } \lambda = \frac{f(b) - f(a)}{g(b) - g(a)} \]
Проверяваме условията на Рол за $H(x)$ в $[a, b]$:
\begin{enumerate}
    \item $H(x)$ е непрекъсната в $[a, b]$ и диференцируема в $(a, b)$ като линейна комбинация от непрекъснати и диференцируеми функции. Нейната производна е $H'(x) = f'(x) - \lambda \cdot g'(x)$.
    \item Проверяваме равенството в краищата:
    \[ H(a) = f(a) - \frac{f(b) - f(a)}{g(b) - g(a)} \cdot g(a) = \frac{f(a)g(b) - f(a)g(a) - f(b)g(a) + f(a)g(a)}{g(b) - g(a)} = \frac{f(a)g(b) - f(b)g(a)}{g(b) - g(a)} \]
    \[ H(b) = f(b) - \frac{f(b) - f(a)}{g(b) - g(a)} \cdot g(b) = \frac{f(b)g(b) - f(b)g(a) - f(b)g(b) + f(a)g(b)}{g(b) - g(a)} = \frac{f(a)g(b) - f(b)g(a)}{g(b) - g(a)} \]
    Следователно $H(a) = H(b)$.
\end{enumerate}

Прилагаме теоремата на Рол за $H(x)$ в интервала $[a, b]$, откъдето следва, че съществува точка $c \in (a, b)$, за която $H'(c) = 0$:
\[ H'(c) = f'(c) - \lambda \cdot g'(c) = 0 \quad \Rightarrow \quad f'(c) = \frac{f(b) - f(a)}{g(b) - g(a)} \cdot g'(c) \]
Тъй като по условие $g'(c) \neq 0$, можем да разделим двете страни на $g'(c)$ и получаваме окончателния вид на теоремата на Коши:
\[ \frac{f(b) - f(a)}{g(b) - g(a)} = \frac{f'(c)}{g'(c)} \]
С това теоремата е напълно доказана.
\end{proof}

\newpage

\section{Формула на Тейлър с остатъчен член във формата на Лагранж}

Формулата на Тейлър ни позволява да апроксимираме достатъчно гладки функции чрез полиноми (многочлени), което е изключително важно за практически пресмятания.

\begin{definition}[Полином на Тейлър]
Нека функцията $f$ притежава производни до ред $n$ включително в точката $x_0$. Полиномът:
\[ T_n(x) = f(x_0) + \frac{f'(x_0)}{1!}(x - x_0) + \frac{f''(x_0)}{2!}(x - x_0)^2 + \dots + \frac{f^{(n)}(x_0)}{n!}(x - x_0)^n = \sum_{k=0}^{n} \frac{f^{(k)}(x_0)}{k!}(x - x_0)^k \]
се нарича \textbf{Полином на Тейлър} от ред $n$ за функцията $f$ в точката $x_0$.
\end{definition}

Ако означим разликата между функцията и нейния полином на Тейлър с $R_n(x)$, получаваме \textbf{Формулата на Тейлър}:
\[ f(x) = T_n(x) + R_n(x) \]
където $R_n(x) = f(x) - T_n(x)$ се нарича \textbf{остатъчен член}. За да бъде формулата полезна, ни е необходима точна оценка за този остатъчен член.

\begin{theorem}[Формула на Тейлър с остатъчен член на Лагранж]
Нека функцията $f$ е непрекъсната заедно с производните си до ред $n$ в затворения интервал с краища $x_0$ и $x$, и съществува $f^{(n+1)}(t)$ във вътрешността на този интервал. Тогава между $x_0$ и $x$ съществува точка $c$, такава че остатъчният член $R_n(x)$ има вида:
\[ R_n(x) = \frac{f^{(n+1)}(c)}{(n+1)!}(x - x_0)^{n+1} \]
\end{theorem}

\begin{proof}
Ще извършим доказателството чрез класическия и академичен метод, показан в университетските лекции, въвеждайки общата помощна функция на \textbf{Шлемилх и Рош}.

Фиксираме точката $x \neq x_0$. Разглеждаме интервала $I$ с краища $x_0$ и $x$. Дефинираме помощна функция $\varphi(t)$ за $t \in I$ по следния начин:
\[ \varphi(t) = f(x) - \sum_{k=0}^{n} \frac{f^{(k)}(t)}{k!}(x - t)^k \]
Разписваме подробно сумата:
\[ \varphi(t) = f(x) - \left[ f(t) + \frac{f'(t)}{1!}(x - t) + \frac{f''(t)}{2!}(x - t)^2 + \dots + \frac{f^{(n)}(t)}{n!}(x - t)^n \right] \]

Нека пресметнем стойностите на функцията $\varphi(t)$ в краищата на интервала ($t = x_0$ и $t = x$):
\begin{enumerate}
    \item За $t = x$: Всички членове от вида $(x - t)^k$ стават нула за $k \ge 1$. Остава само:
    \[ \varphi(x) = f(x) - f(x) = 0 \]
    \item За $t = x_0$: Изразът съвпада точно с дефиницията за остатъчния член $R_n(x)$:
    \[ \varphi(x_0) = f(x) - T_n(x) = R_n(x) \]
\end{enumerate}

Сега диференцираме $\varphi(t)$ по отношение на променливата $t$, като внимаваме, че членовете в сумата са произведения на две функции на $t$, т.е. прилагаме правилото $(u \cdot v)' = u'v + uv'$:
\[ \varphi'(t) = 0 - \left[ f'(t) + \left( \frac{f''(t)}{1!}(x-t) - f'(t) \right) + \left( \frac{f^{(3)}(t)}{2!}(x-t)^2 - \frac{f''(t)}{1!}(x-t) \right) + \dots \right] \]
Забелязваме, че се получава телескопична сума, в която членовете се съкращават по двойки. Единственият член, който не се съкращава от последното диференциране, е:
\[ \varphi'(t) = -\frac{f^{(n+1)}(t)}{n!}(x - t)^n \]

Въвеждаме втора помощна функция $\psi(t) = (x - t)^p$, където $p \ge 1$ е параметър (естествено число). Нейната производна по $t$ е:
\[ \psi'(t) = -p(x - t)^{p-1} \]
За функциите $\varphi(t)$ и $\psi(t)$ са изпълнени условията на \textbf{Теоремата на Коши} в интервала с краища $x_0$ и $x$. Следователно съществува вътрешна точка $c$ между $x_0$ и $x$, такава че:
\[ \frac{\varphi(x) - \varphi(x_0)}{\psi(x) - \psi(x_0)} = \frac{\varphi'(c)}{\psi'(c)} \]

Заместваме известните стойности $\varphi(x) = 0$, $\varphi(x_0) = R_n(x)$, $\psi(x) = 0$ и $\psi(x_0) = (x - x_0)^p$:
\[ \frac{0 - R_n(x)}{0 - (x - x_0)^p} = \frac{-\frac{f^{(n+1)}(c)}{n!}(x - c)^n}{-p(x - c)^{p-1}} \]
\[ \frac{R_n(x)}{(x - x_0)^p} = \frac{f^{(n+1)}(c)}{n! \cdot p} (x - c)^{n - p + 1} \]
Оттук изразяваме общата форма на остатъчния член (Форма на Шлемилх-Рош):
\[ R_n(x) = \frac{f^{(n+1)}(c)}{n! \cdot p} (x - x_0)^p (x - c)^{n - p + 1} \]

За да получим търсената от конспективното изискване \textbf{форма на Лагранж}, избираме стойност на параметъра $p = n + 1$:
\[ R_n(x) = \frac{f^{(n+1)}(c)}{n! \cdot (n + 1)} (x - x_0)^{n+1} (x - c)^{n - (n+1) + 1} \]
Тъй като $n! \cdot (n+1) = (n+1)!$ и $(x - c)^0 = 1$, изразът се опростява до:
\[ R_n(x) = \frac{f^{(n+1)}(c)}{(n+1)!}(x - x_0)^{n+1} \]
С това формулата на Тейлър с остатъчен член във формата на Лагранж е напълно изведена и доказана.
\end{proof}

\end{FlushLeft}

\end{document}